% Part: normal-modal-logic
% Chapter: filtrations
% Section: examples-of-filtrations

\documentclass[../../../include/open-logic-section]{subfiles}

\begin{document}

\olfileid{nml}{fil}{exf}

\olsection{निस्यंदनांची उदाहरणे}

निस्यंदने अस्तित्वात आहेत हे आपण अजून दाखवलेले नाही.
पण प्रत्यक्षात कोणतेही प्रतिमान~$\mModel{M}$ घेतले,
तरी $\mModel{M}$ ची $\Gamma$ मधून मिळणारी अनेक निस्यंदने असतात.
त्यांपैकी दोन विशेष निस्यंदने पाहू: सर्वांत सूक्ष्म
आणि सर्वांत स्थूल. त्याच प्रतिमानाच्या निस्यंदनांचे
प्राप्यता संबंध वेगळे असू शकतात
(\olref[fil]{defn:filtration} मध्ये $W^*$ आणि~$V^*$
थेट निश्चित केले आहेत). सर्वांत सूक्ष्म निस्यंदनात
संबंधित जगांच्या जोड्या शक्य तितक्या कमी असतात,
तर सर्वांत स्थूल निस्यंदनात त्या शक्य तितक्या अधिक असतात.

\begin{defn}
  $\Gamma$ उपसूत्रांखाली बंद असेल, तर प्रतिमान~$\mModel{M}$
  चे \emph{सर्वांत सूक्ष्म} निस्यंदन~$\mModel{M^*}$ पुढीलप्रमाणे
  परिभाषित होते:
  \[
  R^*[u][v] \quad \text{तेव्हा आणि तेव्हाच} \quad \exists u'\in [u] \;
  \exists v' \in [v] : Ru'v'.
  \]
\end{defn}

\begin{prop}\ollabel{prop:finest}
  सर्वांत सूक्ष्म निस्यंदन~$\mModel{M^*}$ हे खरोखर निस्यंदन आहे.
\end{prop}

\begin{proof}
  अशा प्रकारे परिभाषित $R^*$ हा
  \olref[fil]{defn:filtration}\olref[fil]{defn:filtration-R}
  मधील अटी पूर्ण करतो हे तपासायचे आहे. तिन्ही अटी क्रमाने पाहू.

  $Ruv$ असेल, तर $u \in [u]$ आणि $v \in [v]$ असल्याने
  $R^*[u][v]$ सुद्धा. म्हणून
  \olref[fil]{defn:filtration-R1} ही अट पूर्ण होते.

  \iftag{prvBox}{\iftag{probBox}{
      \olref[fil]{defn:filtration-R2} ही अट तपासणे सरावासाठी सोडले आहे.}{
        \olref[fil]{defn:filtration-R2} साठी $\Box!A \in \Gamma$,
        $R^*[u][v]$ आणि $\mSat{M}{\Box!A}[u]$ असे समजा.
        $R^*$ च्या व्याख्येवरून $u' \equiv u$ आणि $v' \equiv v$
        अशी जगे आहेत की $Ru'v'$. $u$ आणि $u'$ हे $\Gamma$
        वरील सत्यतेबाबत एकमत असल्याने
        $\mSat{M}{\Box!A}[u']$; म्हणून $\mSat{M}{!A}[v']$.
        $\Gamma$ उप-!!{formula}s खाली बंद असल्याने
        $v$ आणि $v'$ यांच्यात~$!A$ च्या सत्यतेबाबतही
        एकमत आहे. म्हणून अपेक्षेप्रमाणे $\mSat{M}{!A}[v]$.}}{}

  \iftag{prvDiamond}{\iftag{probDiamond}{
      \olref[fil]{defn:filtration-R3} ही अट तपासणे सरावासाठी सोडले आहे.}{
        \olref[fil]{defn:filtration-R3} तपासण्यासाठी
        $\Diamond!A \in \Gamma$, $R^*[u][v]$ आणि
        $\mSat{M}{!A}[v]$ असे समजा. $R^*$ च्या व्याख्येवरून
        $u' \equiv u$ आणि $v' \equiv v$ अशी जगे आहेत की
        $Ru'v'$. $v$ आणि $v'$ यांच्यात~$\Gamma$ वरील
        सत्यतेबाबत एकमत आहे आणि $\Gamma$ उप-!!{formula}s
        खाली बंद आहे; म्हणून $\mSat{M}{!A}[v']$ सुद्धा.
        त्यामुळे $\mSat{M}{\Diamond !A}[u']$.
        $u$ आणि $u'$ यांच्यातही $\Gamma$ वरील
        सत्यतेबाबत एकमत असल्याने
        $\mSat{M}{\Diamond !A}[u]$.}}{}
\end{proof}

\begin{probtag}{probBox,probDiamond}
  \olref[nml][fil][exf]{prop:finest} ची सिद्धता पूर्ण करा.
\end{probtag}

\begin{defn}
  $\Gamma$ उपसूत्रांखाली बंद असेल, तर प्रतिमान~$\mModel{M}$
  चे \emph{सर्वांत स्थूल} निस्यंदन~$\mModel{M^*}$
  पुढीलप्रमाणे परिभाषित करतो:
  $R^*[u][v]$ तेव्हा आणि तेव्हाच
  \iftag{notprvBox,notprvDiamond}{पुढील अट पूर्ण होते:}{
    पुढील \emph{दोन्ही} अटी पूर्ण होतात:}
  \begin{tagenumerate}{prvBox,prvDiamond}
  \tagitem{prvBox}{\ollabel{defn:coarsest-Box}$\Box!A \in \Gamma$
    आणि $\mSat{M}{\Box!A}[u]$ असेल, तर
    $\mSat{M}{!A}[v]$\iftag{prvDiamond}{;}{.}}{}
  \tagitem{prvDiamond}{\ollabel{defn:coarsest-Diamond}$\Diamond!A
    \in \Gamma$ आणि $\mSat{M}{!A}[v]$ असेल, तर
    $\mSat{M}{\Diamond!A}[u]$.}{}
  \end{tagenumerate}
\end{defn}

\begin{prop}
  सर्वांत स्थूल निस्यंदन~$\mModel{M^*}$ हे खरोखर निस्यंदन आहे.
\end{prop}

\begin{proof}
  $R^*$ च्या व्याख्येनुसार आता फक्त $Ruv$ वरून
  $R^*[u][v]$ मिळते का हे तपासायचे आहे. म्हणून
  $Ruv$ असे समजा. \iftag{prvBox}{$\Box!A \in \Gamma$
    आणि $\mSat{M}{\Box!A}[u]$ असे समजा; मग स्पष्टपणे
    $\mSat{M}{!A}[v]$\iftag{prvDiamond}{, आणि
      \olref{defn:coarsest-Box} ही अट पूर्ण होते.}{.}}{}\iftag{prvDiamond}{
    $\Diamond!A \in \Gamma$ आणि $\mSat{M}{!A}[v]$
    असे समजा. $Ruv$ असल्याने $\mSat{M}{\Diamond!A}[u]$
    \iftag{prvBox}{, आणि \olref{defn:coarsest-Diamond}
      ही अट पूर्ण होते}{}.}{}
\end{proof}

\begin{ex}
  $W = \PosInt$, $Rnm$ तेव्हा आणि तेव्हाच $m = n + 1$,
  आणि $V(p) = \Setabs{2n}{n \in
    \PosInt}$ असे समजा. प्रतिमान~$\mModel{M} = \tuple{W, R, V}$
  हे \olref{fig:ex-filtration} मध्ये दाखवले आहे.
  त्याची जगे $1$, $2$ इत्यादी आहेत. प्रत्येक जगाला
  नेमके एकच जग, म्हणजे त्याचे पुढचे जग, प्राप्य आहे;
  आणि $p$ तेव्हा आणि तेव्हाच सत्य आहे जेव्हा संख्या सम असते.

  \begin{figure}
    \centering
    \begin{tikzpicture}[modal]
      \node[world] (1) [label=below:\mFalse{p}] {$1$};
      \node[world] (2) [label=below:\mTrue{p}, right=of 1] {$2$};
      \node[world] (3) [label=below:\mFalse{p}, right=of 2] {$3$};
      \node[world] (4) [label=below:\mTrue{p}, right=of 3] {$4$};
      \node[phantom] (5) [right=of 4] {};
      \draw[->] (1) to (2);
      \draw[->] (2) to (3);
      \draw[->] (3) to (4);
      \draw[dotted] (4) to (5);
    \end{tikzpicture}

    \begin{tikzpicture}[modal]
      \node[world] (1) [label=below:\mFalse{p}] {$[1]$};
      \node[world] (2) [label=below:\mTrue{p}, right=of 1] {$[2]$};
      \draw[->,bend left] (1) to (2);
      \draw[->,bend left] (2) to (1);

      \node[world] (3) [label=below:\mFalse{p},right=of 2] {$[1]$};
      \node[world] (4) [label=below:\mTrue{p}, right=of 3] {$[2]$};
      \draw[->,bend left] (3) to (4);
      \draw[->,bend left] (4) to (3);
      \draw[->,reflexive above] (4) to (4);
    \end{tikzpicture}
    \caption{एक अनंत प्रतिमान आणि त्याची निस्यंदने.}
    \ollabel{fig:ex-filtration}
  \end{figure}

  आता $\Gamma$ हा~$\Box p \lif p$ च्या सर्व
  उप-!!{formula}s चा संच, म्हणजे
  $\{p, \Box p, \Box p \lif p\}$, असू द्या.
  $p$ सम संख्यांवरच सत्य आहे; $\Box p$ विषम
  संख्यांवरच सत्य आहे. म्हणून $\Box p \lif p$
  हे सम संख्यांवरच सत्य आहे. दुसऱ्या शब्दांत,
  प्रत्येक विषम संख्येवर $\Box p$ सत्य, पण $p$
  आणि $\Box p \lif p$ असत्य आहेत; प्रत्येक सम
  संख्येवर $p$ आणि $\Box p \lif p$ सत्य, पण
  $\Box p$ असत्य आहे. म्हणून
  $W^* = \{ [1], [2] \}$, जेथे
  $[1] = \{1, 3, 5, \dots\}$ आणि
  $[2] = \{2, 4, 6, \dots\}$.
  $2 \in V(p)$ असल्याने $[2] \in V^*(p)$;
  $1 \notin V(p)$ असल्याने $[1] \notin V^*(p)$.
  म्हणून $V^*(p) = \{[2]\}$.

  $W^*$ वर आधारलेल्या कोणत्याही निस्यंदनाच्या
  प्राप्यता संबंधात $\tuple{[1], [2]}, \tuple{[2],[1]}$
  या जोड्या असायलाच हव्यात: $R12$ असल्याने
  \olref[fil]{defn:filtration}\olref[fil]{defn:filtration-R1}
  वरून $R^*[1][2]$; तसेच $R23$ असल्याने
  $R^*[2][3]$, आणि $[3]=[1]$.
  त्यात $\tuple{[1],[1]}$ मात्र असू शकत नाही:
  ती जोडी असेल, तर $R^*[1][1]$ होईल;
  पण $\mSat{M}{\Box p}[1]$ आणि
  $\mSat/{M}{p}[1]$ असल्याने
  \olref[fil]{defn:filtration-R2} शी विरोध येईल.
  $R^*[2][2]$ असावे किंवा नसावे, अशी कोणतीही
  अट नाही. म्हणून~$\mModel{M}$ ची दोन संभाव्य
  निस्यंदने आहेत; त्यांचे प्राप्यता संबंध अनुक्रमे
  \[
  \{\tuple{[1],[2]}, \tuple{[2],[1]}\} \text{आणि}
  \{\tuple{[1],[2]}, \tuple{[2],[1]}, \tuple{[2],[2]}\}.
  \]
  असे आहेत. दोन्ही प्रकरणांत~$[1]$ येथे $p$ आणि
  $\Box p \lif p$ असत्य, तर $\Box p$ सत्य आहे;
  $[2]$ येथे $p$ आणि $\Box p \lif p$ सत्य,
  तर $\Box p$ असत्य आहे.
\end{ex}


\begin{prob}
  पुढील प्रतिमान~$\mModel{M} = \tuple{W, R, V}$ पाहा.
  येथे $W = \Setabs{0\sigma}{\sigma \in \Bin^*}$,
  म्हणजे~$0$ ने सुरू होणाऱ्या $0$ व~$1$ यांच्या
  सांत क्रमांचा संच; $R\sigma\sigma'$ तेव्हा आणि तेव्हाच
  $\sigma' = \sigma 0$ किंवा $\sigma' = \sigma 1$;
  तसेच $V(p) = \{0\} \cup \Setabs{\sigma
    0}{\sigma \in W}$ आणि $V(q) = \Setabs{\sigma 1}{\sigma \in
    W}$. त्याचे चित्र असे आहे:
  \begin{center}
    \begin{tikzpicture}[modal,
        node distance=2cm
      ]

      \node[world] (0) [label={below:\mTrue{p}\\\mFalse{q}},font=\tiny] {$0$};
      \node[world] (00) [label={below:\mTrue{p}\\\mFalse{q}},
        above right=of 0,font=\tiny] {$00$};
      \node[world] (000) [label={below:\mTrue{p}\\\mFalse{q}},
        right=of 00, yshift=.8cm, font=\tiny] {$000$};
      \node[phantom] (0000) [right=of 000, yshift=.5cm] {};
      \node[phantom] (0001) [right=of 000, yshift=-.5cm] {};
      \node[world] (001) [label={below:\mFalse{p}\\\mTrue{q}},
        right=of 00, yshift=-.8cm, font=\tiny] {$001$};
      \node[phantom](0010) [right=of 001, yshift=.5cm] {};
      \node[phantom](0011) [right=of 001, yshift=-.5cm] {};
      \node[world] (01) [label={below:\mFalse{p}\\\mTrue{q}},
        below right=of 0,font=\tiny] {$01$};
      \node[world] (010) [label={below:\mTrue{p}\\\mFalse{q}},
        right=of 01, yshift=.8cm,font=\tiny] {$010$};
      \node[phantom] (0100) [right=of 010, yshift=.5cm] {};
      \node[phantom] (0101) [right=of 010, yshift=-.5cm] {};
      \node[world] (011) [label={below:\mFalse{p}\\\mTrue{q}},
        right=of 01, yshift=-.8cm,font=\tiny] {$011$};
      \node[phantom] (0110) [right=of 011, yshift=.5cm] {};
      \node[phantom] (0111) [right=of 011, yshift=-.5cm] {};

      \draw[->] (0) to (00);
      \draw[->] (0) to (01);
      \draw[->] (00) to (000);
      \draw[dotted] (000) to (0000);
      \draw[dotted] (000) to (0001);
      \draw[->] (00) to (001);
      \draw[dotted] (001) to (0010);
      \draw[dotted] (001) to (0011);
      \draw[->] (01) to (010);
      \draw[dotted] (010) to (0100);
      \draw[dotted] (010) to (0101);
      \draw[->] (01) to (011);
      \draw[dotted] (011) to (0110);
      \draw[dotted] (011) to (0111);
    \end{tikzpicture}
  \end{center}
  प्रत्येक~$w$ साठी
  $\mSat/{M}{\Box(p \lor q) \lif (\Box p \lor \Box q)}[w]$
  असे आहे.

  $\Gamma$ हा~$\Box(p \lor q) \lif
  (\Box p \lor \Box q)$ च्या उप-!!{formula}s चा संच
  असू द्या. $W^*$ आणि~$V^*$ काय असतील?
  $\mModel{M}$ च्या सर्वांत सूक्ष्म निस्यंदनाचा
  प्राप्यता संबंध कोणता असेल? सर्वांत स्थूल निस्यंदनाचा कोणता?
\end{prob}

\end{document}
